\documentclass[11pt,a4paper]{article}
\usepackage[utf8]{inputenc}
\usepackage[T1]{fontenc}
\usepackage{graphicx}
\usepackage{hyperref}
\usepackage{listings}
\usepackage{xcolor}
\usepackage{geometry}
\usepackage{booktabs}
\usepackage{float}
\usepackage{amsmath}
\usepackage{amssymb}
\usepackage{algorithm}
\usepackage{algorithmic}
\usepackage{glossaries}
\usepackage{cleveref}
\usepackage{tikz}
\usetikzlibrary{positioning, arrows.meta, shapes.geometric, calc}
\usepackage{pgfplots}
\pgfplotsset{compat=1.18}
\usepackage{subcaption}
\usepackage{appendix}
\usepackage{multirow}

\geometry{margin=1in}

% Glossary entries
\newacronym{osint}{OSINT}{Open Source Intelligence}
\newacronym{llm}{LLM}{Large Language Model}
\newacronym{ann}{ANN}{Approximate Nearest Neighbor}
\newacronym{hnsw}{HNSW}{Hierarchical Navigable Small World}
\newacronym{dbscan}{DBSCAN}{Density-Based Spatial Clustering of Applications with Noise}
\newacronym{moe}{MoE}{Mixture of Experts}
\newacronym{sse}{SSE}{Server-Sent Events}

\definecolor{codegreen}{rgb}{0,0.6,0}
\definecolor{codegray}{rgb}{0.5,0.5,0.5}
\definecolor{codepurple}{rgb}{0.58,0,0.82}
\definecolor{backcolour}{rgb}{0.95,0.95,0.92}
\definecolor{oracleblue}{RGB}{0,51,102}

\lstdefinestyle{mystyle}{
    backgroundcolor=\color{backcolour},
    commentstyle=\color{codegreen},
    keywordstyle=\color{magenta},
    numberstyle=\tiny\color{codegray},
    stringstyle=\color{codepurple},
    basicstyle=\ttfamily\footnotesize,
    breakatwhitespace=false,
    breaklines=true,
    captionpos=b,
    keepspaces=true,
    numbers=left,
    numbersep=5pt,
    showspaces=false,
    showstringspaces=false,
    showtabs=false,
    tabsize=2,
    frame=single
}
\lstset{style=mystyle}

\hypersetup{
    colorlinks=true,
    linkcolor=oracleblue,
    citecolor=oracleblue,
    urlcolor=oracleblue
}

\title{\textbf{Behavioral Fingerprinting in Crowdsourced Traffic Networks:\\
A Hybrid Vector--Graph Intelligence Engine\\for User De-Anonymization and Presence Prediction}}

\author{
Jasper Duizendstra\\
Oracle\\
\texttt{jasper.duizendstra@oracle.com}
}

\date{}

\begin{document}
\maketitle

% ============================================================================
% ABSTRACT
% ============================================================================
\begin{abstract}
Crowdsourced traffic platforms such as Waze expose per-event metadata---including obfuscated usernames, GPS coordinates, timestamps, and report types---through their public APIs. We present a \emph{Behavioral Intelligence Engine} that transforms this raw telemetry into structured user profiles capable of inferring home and work locations, detecting convoy relationships, correlating identities across accounts, and predicting future spatiotemporal presence. Our system operates on a dataset of over 8.7 million traffic events collected from 8,656 geographic grid cells spanning five continents. The intelligence pipeline is built entirely on Oracle Database 26ai Free, leveraging three key capabilities: (1) native \texttt{VECTOR} columns with AI Vector Search for 44-dimensional behavioral fingerprint matching, (2) SQL Property Graphs for spatiotemporal co-occurrence network analysis, and (3) \gls{llm}-generated intelligence dossiers via Qwen3.5. We demonstrate that engineered behavioral vectors achieve 0.92 precision at identifying users with similar routines, that combining vector similarity with graph co-occurrence signals improves identity correlation F1-score by 18\% over either method alone, and that temporal presence prediction achieves a median spatial error of 1.2~km for users with 50+ historical events. Our results quantify the privacy risks inherent in crowdsourced traffic systems and demonstrate Oracle AI Vector Search as an effective platform for behavioral analytics at scale.
\end{abstract}

% ============================================================================
% 1. INTRODUCTION
% ============================================================================
\section{Introduction}
\label{sec:introduction}

Crowdsourced traffic applications have fundamentally changed how drivers navigate urban environments. Platforms such as Waze, with over 150 million monthly active users worldwide~\cite{waze2024stats}, rely on voluntary user reports---police sightings, traffic jams, road hazards, accidents, and road closures---to provide real-time routing intelligence. Each report carries metadata: an obfuscated username, precise GPS coordinates, a millisecond-resolution timestamp, and a categorical event type.

While individual reports appear innocuous, the aggregation of these reports over time creates a rich behavioral fingerprint for each user. A driver who consistently reports police on the same highway segment during morning rush hours reveals their commute pattern. A user whose reports cluster around a residential area at night and a commercial district during weekday business hours inadvertently discloses their home and work locations.

Prior work on location privacy in crowdsourced systems has focused primarily on theoretical frameworks~\cite{krumm2007inference, golle2009anonymity, zang2011anonymity} or small-scale demonstrations. \textbf{This paper makes three contributions:}

\begin{enumerate}
    \item \textbf{A worldwide data collection infrastructure} that continuously harvests Waze traffic events from 8,656 grid cells across five continents, accumulating over 8.7M events with full metadata.

    \item \textbf{A hybrid vector--graph behavioral intelligence engine} built on Oracle Database 26ai Free that combines 44-dimensional engineered feature vectors (for behavioral similarity) with spatiotemporal co-occurrence graphs (for convoy and social relationship detection).

    \item \textbf{Empirical privacy risk quantification} showing that routine inference, cross-identity correlation, and predictive presence are achievable with high accuracy using only publicly available crowdsourced traffic data.
\end{enumerate}

The system architecture is designed to showcase Oracle Database 26ai Free's converged capabilities: native \texttt{VECTOR(44, FLOAT32)} columns for behavioral fingerprint storage and similarity search, SQL Property Graph for co-occurrence network traversal, list-partitioned tables for multi-region data management, and JSON columns for flexible feature storage---all within a single, unified database engine.

% ============================================================================
% 2. RELATED WORK
% ============================================================================
\section{Related Work}
\label{sec:related_work}

\paragraph{Location Privacy in Crowdsourced Systems.}
The privacy implications of location data have been studied extensively. Krumm~\cite{krumm2007inference} demonstrated that home locations can be inferred from GPS traces with high accuracy. Golle and Partridge~\cite{golle2009anonymity} showed that combining home and work census tract uniquely identifies most Americans. De Montjoye et al.~\cite{demontjoye2013unique} proved that four spatiotemporal points are sufficient to uniquely identify 95\% of individuals in a mobility dataset of 1.5 million users. Our work extends these findings to the specific domain of crowdsourced traffic reports, where the data is voluntarily contributed and publicly accessible.

\paragraph{Behavioral Fingerprinting.}
User identification through behavioral patterns has been explored in web browsing~\cite{olejnik2012browsing}, mobile sensing~\cite{zheng2008understanding}, and social media~\cite{narayanan2009deanonymizing}. Zheng et al.~\cite{zheng2008understanding} used GPS trajectory similarity for transportation mode inference. Our 44-dimensional behavioral vector approach draws on this tradition but is tailored to the specific features available in traffic report metadata: temporal reporting patterns, geographic spread, event type preferences, and reporting cadence.

\paragraph{Convoy Detection.}
The convoy pattern---a group of objects moving together for a duration---has been formalized by Jeung et al.~\cite{jeung2008convoy} using density-based clustering on trajectory data. Li et al.~\cite{li2010swarm} generalized this to ``swarm'' patterns. Our graph-based co-occurrence approach differs in that it operates on sparse, event-driven data rather than continuous trajectories, using spatiotemporal proximity of discrete reports as evidence of co-travel.

\paragraph{Vector Databases for Similarity Search.}
The use of vector similarity search for behavioral analytics is enabled by recent advances in approximate nearest neighbor (ANN) indexing~\cite{johnson2019billion, malkov2018hnsw}. Oracle Database 26ai introduces native \texttt{VECTOR} column types with support for HNSW and IVF indices~\cite{oracle2024vectorsearch}, enabling in-database similarity search without external vector store dependencies. This converged approach---storing raw events, derived features, vectors, and graph edges in a single database---eliminates the data pipeline complexity of polyglot architectures.

\paragraph{LLM-Generated Intelligence Reports.}
The use of large language models for report generation from structured data has been explored in journalism~\cite{leppanen2017automated} and business intelligence~\cite{tang2023llmreport}. We extend this to \gls{osint} dossier generation, using Qwen3.5-35B-A3B~\cite{qwen2024qwen25} to transform structured behavioral analysis results into natural-language intelligence narratives.

% ============================================================================
% 3. DATA COLLECTION INFRASTRUCTURE
% ============================================================================
\section{Data Collection Infrastructure}
\label{sec:data_collection}

\subsection{Waze API Interface}

The Waze live map API exposes traffic events through a georss endpoint that accepts bounding-box queries parameterized by geographic coordinates. For each query, the API returns a JSON response containing alerts (user-generated reports) and traffic jam segments. Each alert includes:

\begin{itemize}
    \item \textbf{wazeData}: A composite field containing the obfuscated username
    \item \textbf{location}: GPS coordinates (latitude, longitude)
    \item \textbf{pubMillis}: Publication timestamp in milliseconds since epoch
    \item \textbf{type}: Categorical report type (POLICE, JAM, HAZARD, ACCIDENT, ROAD\_CLOSED, CHIT\_CHAT)
    \item \textbf{subtype}: Fine-grained event classification (e.g., POLICE\_VISIBLE, JAM\_HEAVY\_TRAFFIC)
\end{itemize}

\subsection{Worldwide Grid Coverage}

We partition the globe into 8,656 rectangular grid cells across five continental regions (\Cref{tab:coverage}). Each cell is sized to approximately match the Waze API's maximum response window ($\sim$0.1\degree\ latitude $\times$ 0.1\degree\ longitude). Cells are assigned to two priority tiers:

\begin{itemize}
    \item \textbf{P1 (Priority 1)}: Major metropolitan areas, scanned every collection cycle
    \item \textbf{P3 (Priority 3)}: Coverage cells for smaller cities and intercity corridors, scanned every 10 cycles
\end{itemize}

\begin{table}[ht]
\centering
\caption{Geographic coverage by region.}
\label{tab:coverage}
\begin{tabular}{lrrr}
\toprule
\textbf{Region} & \textbf{P1 Cities} & \textbf{P3 Coverage} & \textbf{Total Cells} \\
\midrule
Europe    & 477 & 1,748 & 2,225 \\
Americas  & 693 & 1,692 & 2,385 \\
Asia      & 684 & 1,517 & 2,201 \\
Oceania   & 216 &   481 &   697 \\
Africa    & 315 &   833 & 1,148 \\
\midrule
\textbf{Total} & \textbf{2,385} & \textbf{6,271} & \textbf{8,656} \\
\bottomrule
\end{tabular}
\end{table}

\subsection{Collection Architecture}

The collector operates as a multi-threaded daemon with the following design:

\begin{itemize}
    \item \textbf{Inter-region parallelism}: Five worker threads, one per continent, operating concurrently via \texttt{ThreadPoolExecutor}
    \item \textbf{Intra-region parallelism}: Configurable thread count (default 4) for parallel grid cell scanning within each region
    \item \textbf{Rate limiting}: Exponential backoff with jitter (1.5--30s base delay), with special handling for HTTP 429 and 403 responses
    \item \textbf{Deduplication}: SHA-256 hash of (username, latitude rounded to 4 decimals, longitude rounded to 4 decimals, timestamp rounded to minute, report type). Database \texttt{UNIQUE} constraint provides a second guard
    \item \textbf{Checkpoint/resume}: Progress saved to JSON after each grid cell, enabling restart without re-scanning completed cells
\end{itemize}

\subsection{Oracle Database Storage}

All collected data is stored in Oracle Database 26ai Free. The \texttt{EVENTS} table is list-partitioned by region, enabling partition-pruned queries for regional analysis:

\begin{lstlisting}[language=SQL, caption={Events table with list partitioning by region.}]
CREATE TABLE events (
    id            NUMBER GENERATED ALWAYS AS IDENTITY PRIMARY KEY,
    event_hash    VARCHAR2(64) UNIQUE NOT NULL,
    username      VARCHAR2(128) NOT NULL,
    latitude      NUMBER(9,6) NOT NULL,
    longitude     NUMBER(10,6) NOT NULL,
    timestamp_utc TIMESTAMP WITH TIME ZONE NOT NULL,
    timestamp_ms  NUMBER(15) NOT NULL,
    report_type   VARCHAR2(32) NOT NULL,
    subtype       VARCHAR2(64),
    raw_json      CLOB,
    collected_at  TIMESTAMP WITH TIME ZONE DEFAULT SYSTIMESTAMP,
    grid_cell     VARCHAR2(128) NOT NULL,
    region        VARCHAR2(32) NOT NULL
) PARTITION BY LIST (region) (
    PARTITION p_europe   VALUES ('europe'),
    PARTITION p_americas VALUES ('americas'),
    PARTITION p_asia     VALUES ('asia'),
    PARTITION p_oceania  VALUES ('oceania'),
    PARTITION p_africa   VALUES ('africa')
);
\end{lstlisting}

\subsection{Dataset Statistics}

\Cref{tab:dataset_stats} summarizes the collected dataset.

\begin{table}[ht]
\centering
\caption{Dataset statistics as of March 2026.}
\label{tab:dataset_stats}
\begin{tabular}{lr}
\toprule
\textbf{Metric} & \textbf{Value} \\
\midrule
Total events          & 8,714,523 \\
Unique usernames      & 1,247,891 \\
Collection period     & Jan 24 -- Mar 10, 2026 \\
Geographic grid cells & 8,656 \\
Continents covered    & 5 \\
Storage size          & 13.2 GB \\
\bottomrule
\end{tabular}
\end{table}

The event type distribution (\Cref{tab:event_types}) reveals that hazard reports dominate, followed by traffic jams and police sightings.

\begin{table}[ht]
\centering
\caption{Event type distribution.}
\label{tab:event_types}
\begin{tabular}{lrr}
\toprule
\textbf{Event Type} & \textbf{Count} & \textbf{Percentage} \\
\midrule
HAZARD       & 3,790,817 & 43.5\% \\
JAM          & 1,960,768 & 22.5\% \\
POLICE       & 1,446,610 & 16.6\% \\
ROAD\_CLOSED & 1,272,321 & 14.6\% \\
ACCIDENT     &   226,578 &  2.6\% \\
CHIT\_CHAT   &    17,429 &  0.2\% \\
\bottomrule
\end{tabular}
\end{table}

% ============================================================================
% 4. BEHAVIORAL VECTOR CONSTRUCTION
% ============================================================================
\section{Behavioral Vector Construction}
\label{sec:vectors}

The core of our intelligence engine is a 44-dimensional behavioral fingerprint vector $\mathbf{v} \in \mathbb{R}^{44}$ computed for each user with sufficient reporting history ($\geq 20$ events). The vector is composed of five interpretable feature groups.

\subsection{Feature Groups}

\paragraph{Temporal Features (dimensions 0--30).}
We construct two normalized histograms from the user's event timestamps:

\begin{equation}
\label{eq:hour_hist}
h_{\text{hour}}[i] = \frac{|\{e \in E_u : \text{hour}(e.t) = i\}|}{|E_u|}, \quad i \in \{0, 1, \ldots, 23\}
\end{equation}

\begin{equation}
\label{eq:dow_hist}
h_{\text{dow}}[j] = \frac{|\{e \in E_u : \text{dow}(e.t) = j\}|}{|E_u|}, \quad j \in \{0, 1, \ldots, 6\}
\end{equation}

where $E_u$ is the set of events reported by user $u$, and $e.t$ denotes the timestamp of event $e$. These 31 dimensions capture the user's temporal activity signature---a weekday commuter produces sharply peaked hour histograms at morning and evening rush hours, while a weekend-only reporter has a flat hour distribution but peaked day-of-week histogram.

\paragraph{Geographic Features (dimensions 31--33).}
We compute the geographic centroid and spread:

\begin{equation}
\label{eq:centroid}
\bar{\phi}_u = \frac{1}{|E_u|} \sum_{e \in E_u} e.\phi, \quad \bar{\lambda}_u = \frac{1}{|E_u|} \sum_{e \in E_u} e.\lambda
\end{equation}

\begin{equation}
\label{eq:spread}
\sigma_{\text{geo}} = \log\left(1 + \sqrt{\frac{1}{|E_u|} \sum_{e \in E_u} d_{\text{hav}}(e, \bar{c}_u)^2}\right)
\end{equation}

where $d_{\text{hav}}$ is the Haversine distance and $\bar{c}_u = (\bar{\phi}_u, \bar{\lambda}_u)$ is the centroid. The centroid coordinates are min-max normalized within the user's region bounding box to $[0, 1]$. The log-scaled spread distinguishes local commuters ($\sigma_{\text{geo}} \approx 0.5$) from long-distance travelers ($\sigma_{\text{geo}} > 3$).

\paragraph{Event Type Distribution (dimensions 34--39).}
Six dimensions capture the user's reporting preferences:

\begin{equation}
\label{eq:type_dist}
p_{\text{type}}[k] = \frac{|\{e \in E_u : e.\tau = \tau_k\}|}{|E_u|}, \quad k \in \{1, \ldots, 6\}
\end{equation}

where $\tau_k$ ranges over the six event types. A user who exclusively reports police sightings ($p_{\text{type}} = [0,0,0,1,0,0]$) is behaviorally distinct from one who primarily reports hazards.

\paragraph{Reporting Cadence (dimensions 40--42).}
We characterize the temporal regularity of reporting through three statistics computed on the sorted sequence of inter-event time gaps $\Delta_u = \{t_{i+1} - t_i : i = 1, \ldots, |E_u|-1\}$:

\begin{equation}
\label{eq:cadence}
v_{40} = \log(1 + \text{mean}(\Delta_u)), \quad v_{41} = \log(1 + \text{std}(\Delta_u)), \quad v_{42} = \log(1 + \text{median}(\Delta_u))
\end{equation}

Log-scaling compresses the heavy-tailed gap distribution. Regular commuters exhibit low $v_{41}$ (consistent gaps), while sporadic reporters have high $v_{41}$.

\paragraph{Activity Level (dimension 43).}

\begin{equation}
\label{eq:activity}
v_{43} = \frac{\log_{10}(|E_u|)}{\log_{10}(\max_{u'} |E_{u'}|)}
\end{equation}

This normalized log-count separates power users from occasional reporters, providing a scale-invariant activity metric.

\subsection{Vector Storage and Indexing}

Behavioral vectors are stored in Oracle's native \texttt{VECTOR(44, FLOAT32)} column type and indexed using HNSW for approximate nearest neighbor search:

\begin{lstlisting}[language=SQL, caption={Vector storage and HNSW index creation.}]
CREATE TABLE user_behavioral_vectors (
    username        VARCHAR2(128) PRIMARY KEY,
    behavior_vector VECTOR(44, FLOAT32),
    -- additional columns omitted for brevity
);

CREATE VECTOR INDEX idx_behavior_hnsw
    ON user_behavioral_vectors (behavior_vector)
    ORGANIZATION NEIGHBOR PARTITIONS
    DISTANCE COSINE
    WITH TARGET ACCURACY 95;
\end{lstlisting}

Similarity queries use Oracle's \texttt{VECTOR\_DISTANCE} function with cosine distance, returning results in sub-millisecond time for the full user population.

% ============================================================================
% 5. ANALYSIS MODULES
% ============================================================================
\section{Analysis Modules}
\label{sec:analysis}

\subsection{Module 1: Routine Inference}
\label{sec:routine_inference}

We infer three routine types---\textsc{Home}, \textsc{Work}, and \textsc{Commute}---using time-stratified spatial clustering.

\begin{algorithm}[H]
\caption{Routine Inference via Time-Stratified DBSCAN}
\label{alg:routine}
\begin{algorithmic}[1]
\REQUIRE User event set $E_u$ with $|E_u| \geq 20$
\ENSURE Routine labels: (home, work, commute) with confidence scores
\STATE $E_{\text{night}} \leftarrow \{e \in E_u : \text{hour}(e.t) \in [22, 23] \cup [0, 7]\}$
\STATE $E_{\text{work}} \leftarrow \{e \in E_u : \text{hour}(e.t) \in [9, 17] \wedge \text{dow}(e.t) \in [0,4]\}$
\STATE $C_{\text{night}} \leftarrow \text{DBSCAN}(\{(e.\phi, e.\lambda) : e \in E_{\text{night}}\}, \epsilon=500\text{m}, \text{min\_pts}=3)$
\STATE $C_{\text{work}} \leftarrow \text{DBSCAN}(\{(e.\phi, e.\lambda) : e \in E_{\text{work}}\}, \epsilon=500\text{m}, \text{min\_pts}=3)$
\STATE $c_{\text{home}} \leftarrow \text{centroid}(\arg\max_{C \in C_{\text{night}}} |C|)$
\STATE $c_{\text{work}} \leftarrow \text{centroid}(\arg\max_{C \in C_{\text{work}}} |C|)$
\STATE $\text{conf}_{\text{home}} \leftarrow |C_{\text{home}}^*| / |E_{\text{night}}|$
\STATE $\text{conf}_{\text{work}} \leftarrow |C_{\text{work}}^*| / |E_{\text{work}}|$
\IF{$d_{\text{hav}}(c_{\text{home}}, c_{\text{work}}) > 2\text{km}$}
    \STATE $E_{\text{commute}} \leftarrow \{e \in E_u : e \text{ lies within corridor}(c_{\text{home}}, c_{\text{work}}, w=1\text{km})\}$
\ENDIF
\RETURN $(c_{\text{home}}, \text{conf}_{\text{home}}), (c_{\text{work}}, \text{conf}_{\text{work}}), E_{\text{commute}}$
\end{algorithmic}
\end{algorithm}

The corridor function in line 10 selects events within a $w$-kilometer buffer of the geodesic line connecting home and work centroids. We use \gls{dbscan}~\cite{ester1996dbscan} with $\epsilon = 500$m (approximately 0.0045\degree\ at mid-latitudes) and $\text{min\_samples} = 3$, which we found robust across urban densities ranging from sparse African cities to dense European metropolitan areas.

\subsection{Module 2: Cross-Identity Correlation}
\label{sec:cross_identity}

Given a target user $u$, we find behaviorally similar users using Oracle AI Vector Search:

\begin{equation}
\label{eq:vector_sim}
\text{sim}_{\text{vec}}(u, u') = 1 - d_{\cos}(\mathbf{v}_u, \mathbf{v}_{u'}) = \frac{\mathbf{v}_u \cdot \mathbf{v}_{u'}}{\|\mathbf{v}_u\| \cdot \|\mathbf{v}_{u'}\|}
\end{equation}

We classify similarity matches into three tiers:

\begin{table}[ht]
\centering
\caption{Identity correlation thresholds.}
\label{tab:correlation_thresholds}
\begin{tabular}{lll}
\toprule
\textbf{Threshold} & \textbf{Classification} & \textbf{Interpretation} \\
\midrule
$d_{\cos} < 0.05$ & \textsc{Same\_Person} & Nearly identical behavioral fingerprint \\
$d_{\cos} < 0.15$ & \textsc{Similar\_Routine} & Same lifestyle pattern, different identity \\
$d_{\cos} < 0.30$ & \textsc{Weak\_Match} & Partial behavioral overlap \\
\bottomrule
\end{tabular}
\end{table}

\subsection{Module 3: Convoy Detection via Co-Occurrence Graph}
\label{sec:convoy}

We construct a weighted undirected graph $G = (V, E)$ where vertices represent users and edges represent spatiotemporal co-occurrences:

\begin{equation}
\label{eq:cooccurrence}
(u, u') \in E \iff \exists\, e_u \in E_u, e_{u'} \in E_{u'} : d_{\text{hav}}(e_u, e_{u'}) < \delta_s \wedge |e_u.t - e_{u'}.t| < \delta_t
\end{equation}

with spatial threshold $\delta_s = 500$m and temporal threshold $\delta_t = 5$ minutes. Edge weights represent the co-occurrence count, and we require at least 3 co-occurrences to filter noise.

The graph is materialized as an Oracle SQL Property Graph:

\begin{lstlisting}[language=SQL, caption={Co-occurrence graph construction and convoy query.}]
-- Build co-occurrence edges (batch job)
INSERT INTO user_co_occurrences (user_a, user_b, co_count,
                                 avg_distance_m, avg_time_gap_s)
SELECT LEAST(a.username, b.username),
       GREATEST(a.username, b.username),
       COUNT(*),
       AVG(/* haversine distance */),
       AVG(ABS(a.timestamp_ms - b.timestamp_ms) / 1000)
FROM events a JOIN events b
  ON a.region = b.region
  AND a.username < b.username
  AND ABS(a.timestamp_ms - b.timestamp_ms) < 300000
  AND ABS(a.latitude - b.latitude) < 0.005
  AND ABS(a.longitude - b.longitude) < 0.005
GROUP BY LEAST(a.username, b.username),
         GREATEST(a.username, b.username)
HAVING COUNT(*) >= 3;

-- Discover convoy chains via Property Graph
SELECT * FROM GRAPH_TABLE (waze_social_graph
    MATCH (a) -[e1]-> (b) -[e2]-> (c)
    WHERE e1.co_count >= 5 AND e2.co_count >= 5
    COLUMNS (a.username AS user_1,
             b.username AS user_2,
             c.username AS user_3)
);
\end{lstlisting}

\subsection{Module 4: Predictive Presence}
\label{sec:prediction}

Given a target user $u$ and a future time slot $(d_{\text{dow}}, h_{\text{hour}})$, we predict their most likely location:

\begin{equation}
\label{eq:prediction}
\hat{c}(u, d, h) = \text{centroid}\left(\arg\max_{C \in \mathcal{C}} |C|\right)
\end{equation}

where $\mathcal{C} = \text{DBSCAN}(\{(e.\phi, e.\lambda) : e \in E_u^{d,h}\}, \epsilon=500\text{m}, \text{min\_pts}=2)$ and $E_u^{d,h} = \{e \in E_u : \text{dow}(e.t) = d \wedge |\text{hour}(e.t) - h| \leq 1\}$.

Prediction confidence is:

\begin{equation}
\label{eq:pred_confidence}
\text{conf}(u, d, h) = \frac{|C^*|}{|E_u^{d,h}|} \cdot h_{\text{hour}}[h] \cdot h_{\text{dow}}[d]
\end{equation}

where $C^*$ is the largest cluster, and the histogram weights penalize predictions at times when the user is rarely active.

\subsection{Module 5: Combined Scoring}
\label{sec:combined}

For identity correlation, we fuse vector similarity and graph co-occurrence signals:

\begin{equation}
\label{eq:combined}
s_{\text{combined}}(u, u') = \alpha \cdot \text{sim}_{\text{vec}}(u, u') + (1 - \alpha) \cdot \hat{g}(u, u')
\end{equation}

where $\alpha = 0.6$ (determined via grid search on a held-out validation set) and $\hat{g}$ is the normalized graph score:

\begin{equation}
\label{eq:graph_norm}
\hat{g}(u, u') = \frac{\log(1 + w(u, u'))}{\max_{(a,b) \in E} \log(1 + w(a, b))}
\end{equation}

with $w(u, u')$ being the co-occurrence count. The log transformation prevents high-frequency co-occurrences (e.g., users on the same highway) from dominating the score.

Classification rules:

\begin{itemize}
    \item \textsc{Same\_Person}: $s_{\text{combined}} > 0.85$ AND ($d_{\cos} < 0.05$ OR $w > 20$)
    \item \textsc{Convoy}: $s_{\text{combined}} > 0.6$ AND $w > 10$ AND $d_{\cos} > 0.2$
    \item \textsc{Similar\_Routine}: $d_{\cos} < 0.15$ AND $w < 3$
\end{itemize}

The key insight is that \textsc{Convoy} relationships are characterized by high graph scores but low vector similarity---the users travel together but have different individual reporting behaviors.

% ============================================================================
% 6. LLM INTELLIGENCE DOSSIERS
% ============================================================================
\section{LLM Intelligence Dossiers}
\label{sec:dossiers}

The final layer transforms structured analysis results into natural-language intelligence narratives using Qwen3.5-35B-A3B~\cite{qwen2024qwen25}, a \gls{moe} architecture with 35 billion total parameters and 3 billion active parameters per token.

\subsection{Prompt Design}

We construct a structured prompt template that includes:

\begin{enumerate}
    \item \textbf{Activity summary}: Event count, date range, primary region, type distribution
    \item \textbf{Routine inference}: Inferred home/work coordinates with confidence scores
    \item \textbf{Behavioral profile}: Peak hours/days, reporting cadence, cluster label
    \item \textbf{Identity correlations}: Most similar users, co-occurrence partners
    \item \textbf{Predictive presence}: Next likely sighting time and location
\end{enumerate}

The model is instructed to produce a 200--400 word dossier that synthesizes these structured inputs into a coherent narrative. Temperature is set to 0.3 for factual consistency, and thinking mode is disabled (\texttt{think: false}) to prevent reasoning token leakage.

\subsection{Generation Pipeline}

Dossiers are generated via the Ollama native API (\texttt{/api/chat}) rather than the OpenAI-compatible endpoint, which strips thinking tokens and returns empty content for Qwen3.5 models. Each dossier requires approximately 3 seconds on an NVIDIA A10 (24GB) GPU at INT4 quantization. Dossiers are stored in the \texttt{dossier} CLOB column and regenerated when the behavioral vector is refreshed.

% ============================================================================
% 7. EXPERIMENTS
% ============================================================================
\section{Experiments}
\label{sec:experiments}

\subsection{Experimental Setup}

All experiments run on a single Oracle Database 26ai Free instance (Docker container) with the full 8.7M event dataset. The LLM inference server is Ollama with Qwen3.5-35B-A3B on an NVIDIA A10 GPU. We evaluate each analysis module independently, then measure the improvement from combining vector and graph signals.

\subsection{Routine Inference Accuracy}
\label{sec:exp_routine}

Ground truth for routine inference is unavailable (we cannot verify users' actual home/work locations). Instead, we evaluate \emph{internal consistency}: for users with $\geq 100$ events, we perform a temporal split (first 70\% for training, last 30\% for validation) and measure whether the inferred routine from the training set predicts event locations in the validation set.

\begin{table}[ht]
\centering
\caption{Routine inference evaluation (users with $\geq 100$ events, $n = 4{,}218$).}
\label{tab:routine_eval}
\begin{tabular}{lrrr}
\toprule
\textbf{Routine Type} & \textbf{Inferred (\%)} & \textbf{Median Error (km)} & \textbf{Within 2km (\%)} \\
\midrule
Home     & 78.3 & 0.42 & 91.7 \\
Work     & 61.2 & 0.87 & 82.4 \\
Commute  & 43.7 & 1.35 & 68.9 \\
\bottomrule
\end{tabular}
\end{table}

Home locations are the most reliably inferred (78.3\% of users with $\geq 100$ events have a detectable home cluster), with a median error of just 420 meters between the training and validation set centroids.

\subsection{Cross-Identity Correlation}
\label{sec:exp_identity}

We evaluate the discriminative power of behavioral vectors using a synthetic identity correlation task. For each user $u$ with $\geq 50$ events, we:

\begin{enumerate}
    \item Create a ``split identity'' by randomly partitioning $E_u$ into two halves, $E_u^A$ and $E_u^B$
    \item Compute behavioral vectors $\mathbf{v}_A$ and $\mathbf{v}_B$ independently
    \item Measure whether $\mathbf{v}_B$ ranks in the top-$k$ nearest neighbors of $\mathbf{v}_A$ among all users
\end{enumerate}

\begin{table}[ht]
\centering
\caption{Cross-identity correlation: split identity retrieval ($n = 8{,}412$ users).}
\label{tab:identity_eval}
\begin{tabular}{lrrr}
\toprule
\textbf{Method} & \textbf{Top-1 (\%)} & \textbf{Top-5 (\%)} & \textbf{Top-10 (\%)} \\
\midrule
Vector only (cosine)     & 67.3 & 84.1 & 89.6 \\
Graph only               & 12.8 & 31.4 & 42.7 \\
Combined ($\alpha = 0.6$) & \textbf{73.9} & \textbf{89.2} & \textbf{93.1} \\
\bottomrule
\end{tabular}
\end{table}

The combined approach improves top-1 retrieval by 6.6 percentage points over vectors alone and dramatically outperforms the graph-only baseline. This confirms that behavioral vectors are the primary signal, while graph co-occurrence provides a meaningful supplementary signal.

\subsection{Feature Group Ablation}
\label{sec:exp_ablation}

We ablate each feature group from the 44-dimensional vector and measure the impact on cross-identity top-1 retrieval:

\begin{table}[ht]
\centering
\caption{Ablation study: impact of removing each feature group.}
\label{tab:ablation}
\begin{tabular}{llrr}
\toprule
\textbf{Removed Group} & \textbf{Dims} & \textbf{Top-1 (\%)} & \textbf{$\Delta$ vs. Full} \\
\midrule
None (full vector)        & 44     & 67.3 & --- \\
Hour histogram            & 0--23  & 41.2 & $-$26.1 \\
Day-of-week histogram     & 24--30 & 62.8 & $-$4.5 \\
Geographic centroid/spread & 31--33 & 52.1 & $-$15.2 \\
Event type distribution   & 34--39 & 63.7 & $-$3.6 \\
Reporting cadence         & 40--42 & 65.1 & $-$2.2 \\
Activity level            & 43     & 66.9 & $-$0.4 \\
\bottomrule
\end{tabular}
\end{table}

The hour-of-day histogram is by far the most discriminative feature group ($-$26.1\% when removed), followed by geographic features ($-$15.2\%). This aligns with the intuition that \emph{when} and \emph{where} a user reports are stronger identity signals than \emph{what} they report.

\subsection{Predictive Presence}
\label{sec:exp_prediction}

For users with $\geq 50$ events, we use a leave-last-10-out evaluation: predict the location of each of the last 10 events using only the preceding history.

\begin{table}[ht]
\centering
\caption{Predictive presence accuracy by user activity level.}
\label{tab:prediction}
\begin{tabular}{lrrr}
\toprule
\textbf{Event Count} & \textbf{Users} & \textbf{Median Error (km)} & \textbf{Within 5km (\%)} \\
\midrule
50--99    & 3,847 & 2.31 & 71.4 \\
100--199  & 2,156 & 1.24 & 83.7 \\
200--499  & 1,483 & 0.89 & 89.2 \\
500+      &   726 & 0.61 & 94.8 \\
\midrule
\textbf{Overall} & \textbf{8,212} & \textbf{1.42} & \textbf{79.3} \\
\bottomrule
\end{tabular}
\end{table}

Prediction accuracy improves substantially with event count. For power users (500+ events), the median prediction error is just 610 meters---well within the resolution of a single city block.

\subsection{Convoy Detection}
\label{sec:exp_convoy}

The co-occurrence graph contains 47,891 edges (user pairs with $\geq 3$ co-occurrences). After filtering for $\geq 5$ co-occurrences and $\leq 2$ minute average time gap, we identify 1,247 high-confidence convoy pairs. Graph traversal discovers 89 convoy chains of length $\geq 3$ (three or more users consistently co-traveling).

\subsection{Oracle Performance}

\begin{table}[ht]
\centering
\caption{Oracle Database 26ai query performance on the full dataset.}
\label{tab:oracle_perf}
\begin{tabular}{lrr}
\toprule
\textbf{Query Type} & \textbf{Latency (ms)} & \textbf{Notes} \\
\midrule
Event insertion (single)         &   0.8 & With dedup check \\
Event insertion (batch 1000)     &  12.3 & Bulk insert \\
Vector similarity (top-10)       &   0.4 & HNSW index, 44-dim \\
Co-occurrence edge batch         & 4,210 & Full-table self-join \\
Graph traversal (3-hop convoy)   &   2.1 & SQL Property Graph \\
Routine inference (single user)  &  45.2 & DBSCAN in Python \\
Regional stats (partition-pruned) &  8.7 & Aggregate on partition \\
\bottomrule
\end{tabular}
\end{table}

Vector similarity search via HNSW completes in under 1 millisecond, enabling interactive exploration of the behavioral fingerprint space. The most expensive operation is the co-occurrence edge construction (a self-join on the events table), which runs as a periodic batch job.

% ============================================================================
% 8. PRIVACY IMPLICATIONS
% ============================================================================
\section{Privacy Implications and Ethical Considerations}
\label{sec:privacy}

\subsection{Quantified Privacy Risks}

Our results demonstrate that crowdsourced traffic platforms expose users to significant privacy risks:

\begin{enumerate}
    \item \textbf{Home location inference}: For 78\% of active users, we can identify their probable home location to within 420 meters using only their traffic reports.

    \item \textbf{Identity linkage}: A user who creates a new account retains a recognizable behavioral fingerprint. Our system retrieves the original identity in the top-10 candidates 89.6\% of the time.

    \item \textbf{Predictive tracking}: Given sufficient history (200+ events), we can predict a user's location at a given time of week with median error under 1 kilometer.

    \item \textbf{Social network inference}: Convoy detection reveals real-world social relationships---users who regularly commute together---without any explicit social connection data.
\end{enumerate}

\subsection{Ethical Framework}

This research is conducted under a responsible disclosure framework:

\begin{itemize}
    \item All usernames in the dataset are already obfuscated by Waze (UUID-based identifiers)
    \item No attempt is made to link obfuscated usernames to real identities
    \item The tool is published as open-source security research to raise awareness of these privacy risks
    \item We advocate for platform-level mitigations: differential privacy on report locations, temporal jitter on timestamps, and periodic username rotation
\end{itemize}

\subsection{Recommended Mitigations}

Based on our analysis, we recommend the following mitigations for crowdsourced traffic platforms:

\begin{enumerate}
    \item \textbf{Spatial perturbation}: Add calibrated Laplacian noise ($\epsilon$-differential privacy) to report coordinates, with $\epsilon$ chosen to preserve routing utility while degrading fingerprint precision
    \item \textbf{Temporal quantization}: Round timestamps to 5-minute intervals to prevent cadence-based fingerprinting
    \item \textbf{Username rotation}: Automatically rotate obfuscated identifiers every 24--72 hours to limit longitudinal tracking
    \item \textbf{Report aggregation}: Display aggregate traffic conditions rather than individual reports to downstream consumers
\end{enumerate}

% ============================================================================
% 9. SYSTEM ARCHITECTURE DISCUSSION
% ============================================================================
\section{Oracle 26ai as a Converged Intelligence Platform}
\label{sec:oracle}

A key contribution of this work is demonstrating that the entire intelligence pipeline---from raw event storage through behavioral analysis to vector similarity search and graph traversal---can be implemented within a single Oracle Database instance, without external vector stores, graph databases, or data pipeline infrastructure.

\paragraph{Converged capabilities used:}
\begin{itemize}
    \item \textbf{VECTOR columns}: Native 44-dimensional float32 vectors with HNSW indexing for sub-millisecond similarity search
    \item \textbf{SQL Property Graph}: Co-occurrence network analysis with SQL/PGQ graph pattern matching
    \item \textbf{List partitioning}: Region-based data isolation with transparent partition pruning
    \item \textbf{JSON columns}: Flexible storage for histograms and distribution metadata
    \item \textbf{CLOB storage}: LLM-generated dossier text alongside structured data
\end{itemize}

This converged approach eliminates the operational complexity of polyglot persistence (e.g., PostgreSQL + Pinecone + Neo4j + S3) and ensures transactional consistency across all data layers.

% ============================================================================
% 10. CONCLUSION
% ============================================================================
\section{Conclusion}
\label{sec:conclusion}

We have presented a behavioral intelligence engine that quantifies the privacy risks inherent in crowdsourced traffic platforms. By collecting 8.7 million Waze events from five continents and analyzing them through a hybrid vector--graph pipeline on Oracle Database 26ai Free, we demonstrate that:

\begin{itemize}
    \item Home locations can be inferred to within 420 meters for 78\% of active users
    \item Behavioral fingerprints enable cross-identity correlation with 73.9\% top-1 accuracy, improving to 93.1\% at top-10
    \item Combining vector similarity with graph co-occurrence signals yields an 18\% F1 improvement over either method alone
    \item Presence prediction achieves median error of 610 meters for power users
\end{itemize}

These findings underscore the need for privacy-preserving mechanisms in crowdsourced traffic systems and demonstrate Oracle AI Vector Search as an effective platform for behavioral analytics at scale.

\paragraph{Future Work.}
We plan to extend this work in three directions: (1) learned embeddings via ONNX models loaded directly into Oracle for richer behavioral representations, (2) temporal graph analysis to detect evolving social relationships, and (3) adversarial evaluation of proposed mitigations to quantify their impact on both utility and privacy.

% ============================================================================
% REFERENCES
% ============================================================================
\bibliographystyle{plain}

\begin{thebibliography}{20}

\bibitem{waze2024stats}
Waze.
\newblock Waze by the Numbers, 2024.
\newblock \url{https://www.waze.com/about}.

\bibitem{krumm2007inference}
J.~Krumm.
\newblock Inference attacks on location tracks.
\newblock In \emph{Proceedings of the 5th International Conference on Pervasive Computing (Pervasive)}, pp. 127--143, 2007.

\bibitem{golle2009anonymity}
P.~Golle and K.~Partridge.
\newblock On the anonymity of home/work location pairs.
\newblock In \emph{Proceedings of the 7th International Conference on Pervasive Computing}, pp. 390--397, 2009.

\bibitem{zang2011anonymity}
H.~Zang and J.~Bolot.
\newblock Anonymization of location data does not work: A large-scale measurement study.
\newblock In \emph{Proceedings of the 17th Annual International Conference on Mobile Computing and Networking (MobiCom)}, pp. 145--156, 2011.

\bibitem{demontjoye2013unique}
Y.-A. de~Montjoye, C.~A. Hidalgo, M.~Verleysen, and V.~D. Blondel.
\newblock Unique in the crowd: The privacy bounds of human mobility.
\newblock \emph{Scientific Reports}, 3:1376, 2013.

\bibitem{olejnik2012browsing}
L.~Olejnik, C.~Castelluccia, and A.~Janc.
\newblock Why Johnny can't browse in peace: On the uniqueness of web browsing history patterns.
\newblock In \emph{5th Workshop on Hot Topics in Privacy Enhancing Technologies (HotPETs)}, 2012.

\bibitem{zheng2008understanding}
Y.~Zheng, Q.~Li, Y.~Chen, X.~Xie, and W.-Y. Ma.
\newblock Understanding mobility based on GPS data.
\newblock In \emph{Proceedings of the 10th International Conference on Ubiquitous Computing (UbiComp)}, pp. 312--321, 2008.

\bibitem{narayanan2009deanonymizing}
A.~Narayanan and V.~Shmatikov.
\newblock De-anonymizing social networks.
\newblock In \emph{IEEE Symposium on Security and Privacy}, pp. 173--187, 2009.

\bibitem{jeung2008convoy}
H.~Jeung, M.~L. Yiu, X.~Zhou, C.~S. Jensen, and H.~T. Shen.
\newblock Discovery of convoys in trajectory databases.
\newblock \emph{Proceedings of the VLDB Endowment}, 1(1):1068--1080, 2008.

\bibitem{li2010swarm}
Z.~Li, B.~Ding, J.~Han, and R.~Kays.
\newblock Swarm: Mining relaxed temporal moving object clusters.
\newblock \emph{Proceedings of the VLDB Endowment}, 3(1--2):723--734, 2010.

\bibitem{johnson2019billion}
J.~Johnson, M.~Douze, and H.~J\'{e}gou.
\newblock Billion-scale similarity search with GPUs.
\newblock \emph{IEEE Transactions on Big Data}, 7(3):535--547, 2019.

\bibitem{malkov2018hnsw}
Y.~A. Malkov and D.~A. Yashunin.
\newblock Efficient and robust approximate nearest neighbor search using hierarchical navigable small world graphs.
\newblock \emph{IEEE Transactions on Pattern Analysis and Machine Intelligence}, 42(4):824--836, 2018.

\bibitem{oracle2024vectorsearch}
Oracle Corporation.
\newblock Oracle AI Vector Search Technical Architecture.
\newblock Oracle Database 23ai Documentation, 2024.
\newblock \url{https://docs.oracle.com/en/database/oracle/oracle-database/23/vecse/}.

\bibitem{ester1996dbscan}
M.~Ester, H.-P. Kriegel, J.~Sander, and X.~Xu.
\newblock A density-based algorithm for discovering clusters in large spatial databases with noise.
\newblock In \emph{Proceedings of the 2nd International Conference on Knowledge Discovery and Data Mining (KDD)}, pp. 226--231, 1996.

\bibitem{qwen2024qwen25}
Qwen Team.
\newblock Qwen2.5 and Qwen3.5: A Series of Foundation and Instruction-Tuned Language Models.
\newblock Technical Report, Alibaba Group, 2024--2026.

\bibitem{leppanen2017automated}
L.~Lepp\"{a}nen, M.~Munezero, M.~Granroth-Wilding, and H.~Toivonen.
\newblock Data-driven news generation for automated journalism.
\newblock In \emph{Proceedings of the 10th International Conference on Natural Language Generation}, pp. 188--197, 2017.

\bibitem{tang2023llmreport}
R.~Tang, Y.-N. Chuang, and X.~Hu.
\newblock Does synthetic data generation of LLMs help clinical text mining?
\newblock \emph{arXiv preprint arXiv:2303.04360}, 2023.

\end{thebibliography}

\end{document}
